\subsection{Preference Optimization, and Who Holds the Pen}
\label{sec:4.2.2}
Reinforcement learning trains an agent by rewarding what a designer wants and penalizing what they do not, and the agent's values are then whatever maximizes the reward. That is the method's strength and the reason it cannot supply a floor: everything the system cares about traces back to a signal somebody wrote and somebody can rewrite.

Writing that signal directly turned out to be the wrong way to get one. What deployed systems use instead is a reward learned from comparisons — show a person two candidate outputs, record which they prefer, fit a model to the preferences, and optimize against the fitted model. Paul Christiano and colleagues introduced the technique in 2017 \autocite{christiano2017deep}; the 2022 instruction-tuning work applied it at the scale of a deployed assistant \autocite{ouyang2022training}. The signal now comes from a population rather than from one author, which is a real improvement in coverage and changes nothing about who commissions the population.

\runin{Two failures that follow from the structure}
The first is reward hacking, where an agent finds something that pays out without doing the intended job — a system built to flag propaganda that learns to flag everything, because indiscriminate flagging scores as well as discrimination and is easier to produce. Optimizing hard against a learned reward model sharpens this rather than softening it, since the proxy is now itself a fitted approximation with its own soft spots. The second is that a reward written for one setting does not carry: an agent trained to cooperate in a game does not thereby cooperate at work, and alignment demonstrated on a benchmark is a claim about the benchmark.

\runin{Removing the reward model does not remove the author}
Direct preference optimization, published by Rafael Rafailov and coauthors in 2023, showed that the separate reward-model fitting step can be dispensed with: the same objective can be reached with a classification loss over the preference pairs themselves \autocite{rafailov2023direct}. It is a substantial simplification, and what it simplifies is the component and not the custodian. Somebody still chose the pairs, hired the raters, wrote the instructions they rate by, and decides when to run the procedure again. Section~\ref{sec:11.5} covers the documented failure of the resulting signal, which is the same finding from the other end: a method that optimizes for approval produces a system bad at being right against the party approving.

\runin{The pen}
Regularization writes a constraint into the objective, penalizing a class of behavior rather than waiting to reward its absence, which is how a recommender is pushed away from the engagement-maximizing equilibrium that produces filter bubbles. It is an improvement and it does not change who holds the pen. A ranked population can be re-sampled and a penalty term can be set to zero, by the party running the training, in an afternoon. Chapter~\ref{sec:7} is what happens to the ranked population specifically, when the disagreement it was collected to represent is compressed into a scalar.

\runin{Exploration is a moral setting wearing a hyperparameter's clothes}
A learner has to exploit what it already knows and explore what it does not. In a treatment planner, exploration means offering a patient something less established and exploitation means declining to look. Safe exploration — searching inside stated constraints rather than checking for violations afterward — is the technique that makes the trade explicit, and section~\ref{sec:6.3.3} takes it up as a safety practice. What matters here is that the balance point is not tuned to performance. It is where a system decides whose interests bear the cost of its learning, and that is the same kind of decision as deciding what the floor covers: made in advance, by somebody, and better made in public.
